\section{An arithmetic criterion for the measure type}
\label{cls:section-criterion}

The criterion in this section is expressed through finite sets of algebraic
units and the coefficient vectors of their minimal polynomials.  Its proof
has two parts.  A sufficiently large cluster of coefficient vectors forces
concentration of the Bernoulli convolution.  Conversely, concentration
produces many short polynomial prefixes, each of which can be completed to
an irreducible polynomial with a root just above the parameter.

For integers $n\geq 1$ and $d\geq 2$, let $R(n,d)$ be the set of real
algebraic units $r\in(1/2,1)$ of degree $dn$ whose minimal polynomial has
all coefficients in $\{-1,0,1\}$ and is irreducible modulo $2$.  An
algebraic unit here means an algebraic integer with algebraic-integer
inverse.  Write $p_r$ for the minimal polynomial with its sign chosen so
that its constant coefficient is $-1$, and set
\[
  \pi_n(p_r)=([t^0]p_r,\ldots,[t^{n-1}]p_r)\in\{-1,0,1\}^n.
\]
For $r\in R(n,d)$, define the finite set
\begin{equation}
  W_n(r;d)=
  \left\{\pi_n(p_q):q\in R(n,d),\ |q-r|\leq 4^{-n}\right\}.
  \label{cls:projection-set}
\end{equation}
This is a set of vectors: repeated projections, including those arising
from conjugate roots of one polynomial, are counted only once.

For a bit word $u=(u_0,\ldots,u_{n-1})\in\{0,1\}^n$, let
\[
  \mathcal O_u=
  \{w\in\mathbb R^n:(2u_i-1)w_i\geq0\text{ for }0\leq i<n\}
\]
be the associated closed orthant.  A vector with a zero coordinate belongs
to both choices of sign in that coordinate.  For $j\geq1$, put
\begin{equation}
  R_*(n,d,j)=
  \left\{r\in R(n,d):
    \#\left\{u\in\{0,1\}^n:
       \#(W_n(r;d)\cap\mathcal O_u)\geq j(2r)^n
    \right\}\geq\frac{2^n}{d}\right\}.
  \label{cls:retained-roots}
\end{equation}
All comparisons in this definition involve algebraic numbers and finite
sets.  In particular, $R(n,d)$ and $R_*(n,d,j)$ can be obtained by finite
enumeration of integer polynomials and exact algebraic comparisons.

\begin{theorem}[Arithmetic criterion]\label{cls:classification}
Define the parameter set
\begin{equation}
  \mathcal C=(0,1/2)\ \cup\
  \left[(1/2,1)\cap
  \bigcup_{d\geq2}\ \bigcap_{j\geq1}\
  \limsup_{n\to\infty}
  \bigcup_{r\in R_*(n,d,j)}[r-4^{-n},r)\right].
  \label{cls:criterion}
\end{equation}
For every $\lambda\in(0,1)$, the measure $\nu_\lambda$ is singular if and
only if $\lambda\in\mathcal C$, and is absolutely continuous if and only
if $\lambda\notin\mathcal C$.  Replacing $\limsup$ by $\liminf$ in
\eqref{cls:criterion} gives the same set.
\end{theorem}

Thus the degree multiplier $d$ may be fixed before the crowding level
$j$ is chosen.  For a singular parameter in $(1/2,1)$, each level $j$
is attained at every sufficiently large $n$, with the threshold in $n$
allowed to depend on $j$.  The predicate uses only polynomial arithmetic,
one-sided approximation, and counts of coefficient vectors.  Determining
membership for a particular parameter can nevertheless be difficult; the
theorem does not supply a finite decision procedure for arbitrary real
inputs.

\subsection{Measure-theoretic preliminaries}
\label{cls:subsection-preliminaries}

We use the bit form
\[
  X_\lambda=\sum_{i=0}^\infty u_i\lambda^i,
  \qquad \mathbb P(u_i=0)=\mathbb P(u_i=1)=\tfrac12,
\]
with law $\mu_\lambda$, and write $L=(1-\lambda)^{-1}$.
The affine identity $Y_\lambda=2X_\lambda-L$ shows that $\mu_\lambda$
and $\nu_\lambda$ have the same measure type.  For a word
$u\in\{0,1\}^n$, put
\begin{equation}
  s_u(\lambda)=\sum_{i=0}^{n-1}u_i\lambda^i.
  \label{cls:prefix-sum}
\end{equation}
Conditioned on this prefix, $X_\lambda$ lies in
$s_u(\lambda)+\lambda^n[0,L]$.

\begin{lemma}[Pure type and the baseline parameters]
\label{cls:pure-type}
For $0<\lambda<1$, the measure $\mu_\lambda$ is either singular or
absolutely continuous.  It is singular for $\lambda<1/2$ and uniform on
$[0,2]$ for $\lambda=1/2$.  Consequently, $\nu_{1/2}$ is uniform on
$[-2,2]$.
\end{lemma}

\begin{proof}
Let $T_i(x)=i+\lambda x$ for $i=0,1$, and let
$\mathcal T\rho=\tfrac12(T_0)_*\rho+\tfrac12(T_1)_*\rho$.
Iteration from any probability measure $\rho$ gives the law of
\[
  \sum_{i=0}^{n-1}u_i\lambda^i+\lambda^n Z,
\]
where $Z$ has law $\rho$ and is independent of the bits.  Since
$\lambda^nZ\to0$ in probability, these laws converge weakly to
$\mu_\lambda$.  Hence $\mu_\lambda$ is the unique invariant probability
of $\mathcal T$.

The operator $\mathcal T$ preserves absolute continuity and singularity.
By uniqueness of the Lebesgue decomposition, the absolutely continuous
and singular parts of its invariant probability are separately invariant.
If both were nonzero, normalizing them would give two distinct invariant
probabilities, a contradiction.

For $\lambda<1/2$, the support is covered, for every $n$, by $2^n$
intervals of length $L\lambda^n$.  Its Lebesgue measure is therefore
zero.  For $\lambda=1/2$, the usual fair binary expansion gives the uniform
law on $[0,2]$.
\end{proof}

We shall also use the following elementary formulation of concentration.
It isolates the measure-theoretic input needed to establish membership in
\eqref{cls:criterion}.

\begin{lemma}[Compact concentration]\label{cls:compact-concentration}
A Borel probability measure $\rho$ on $\mathbb R$ is singular if and
only if, for every $\varepsilon>0$, there is a compact set $K$ with
\[
  \operatorname{Leb}(K)<\varepsilon,
  \qquad \rho(K)>1-\varepsilon.
\]
\end{lemma}

\begin{proof}
If $\rho$ is singular, apply inner regularity to a Borel null set of full
$\rho$-measure.  Conversely, choose such compact sets $K_m$ with
$\varepsilon=2^{-m}$.  The set $\limsup_m K_m$ has Lebesgue measure
zero, since $\sum_m\operatorname{Leb}(K_m)<\infty$, and has
$\rho$-measure one, since every union $\bigcup_{m\geq M}K_m$ has
$\rho$-measure one.
\end{proof}

\subsection{Algebraic clusters imply singularity}
\label{cls:subsection-forward}

Suppose $\lambda\in(1/2,1)$ satisfies the condition in
\eqref{cls:criterion}, and fix a degree multiplier $d$ witnessing it.
For each $j$, there are arbitrarily large $n$ and
$r\in R_*(n,d,j)$ such that
\[
  r-4^{-n}\leq\lambda<r.
\]
If $w\in W_n(r;d)$, choose a polynomial $p_q$ giving this projection.
Its root $q$ satisfies $|q-\lambda|\leq2\cdot4^{-n}$.  On any fixed
compact subinterval of $(0,1)$ containing $\lambda$ in its interior,
\[
  |p_q'(t)|\leq\sum_{k\geq1}k t^{k-1}=\frac1{(1-t)^2}.
\]
The coefficients of index at least $n$ contribute at most
$L\lambda^n$ at $\lambda$.  Thus there is a constant $B=B(\lambda)$,
independent of $n,d,j,w$, such that for all sufficiently large $n$,
\begin{equation}
  \left|\sum_{i=0}^{n-1}w_i\lambda^i\right|\leq B\lambda^n.
  \label{cls:slab}
\end{equation}
Here we used $4^{-n}=o(\lambda^n)$.

If $w\in\mathcal O_u\cap\{-1,0,1\}^n$, then $v=u-w$ belongs to
$\{0,1\}^n$.  Different $w$ give different $v$.  By
\eqref{cls:retained-roots}, at least $2^n/d$ words $u$ consequently have
at least $j(2r)^n\geq j(2\lambda)^n$ neighbors $v$ satisfying
\[
  |s_u(\lambda)-s_v(\lambda)|\leq B\lambda^n.
\]

Partition the line into half-open bins of length $\lambda^n$, and let
$\mathcal I$ be the bins containing at least one of these words $u$,
where a word occupies the bin containing its prefix sum.  For each
$J\in\mathcal I$, choose one such word.  Its neighbors lie in the
$B\lambda^n$-enlargement of $J$.  Each prefix sum can belong to the
enlargements of at most $N_B=2\lceil B\rceil+3$ bins.  Counting
word--bin incidences yields
\[
  \#\mathcal I\,j(2\lambda)^n\leq N_B2^n,
  \qquad
  \#\mathcal I\leq\frac{N_B\lambda^{-n}}j.
\]
Enlarge each bin $J=[a,a+\lambda^n)$ to
$[a,a+(1+L)\lambda^n]$, and denote their union by $E_j$.  The tail bound
following \eqref{cls:prefix-sum} gives
\begin{equation}
  \operatorname{Leb}(E_j)\leq\frac{(1+L)N_B}{j},
  \qquad \mu_\lambda(E_j)\geq\frac1d.
  \label{cls:forward-concentration}
\end{equation}
Absolute continuity of a finite measure makes its mass uniformly small
on sets of sufficiently small Lebesgue measure.  Therefore
\eqref{cls:forward-concentration}, as $j\to\infty$, rules out absolute
continuity.  Lemma~\ref{cls:pure-type} proves singularity.

\section{Constructing the algebraic approximants}
\label{cls:section-converse}

We now prove the converse in Theorem~\ref{cls:classification}.  Throughout
this section, $\lambda\in(1/2,1)$ is fixed, and we abbreviate
$\mu=\mu_\lambda$, $X=X_\lambda$, and $s_u=s_u(\lambda)$.  The first
step is to turn concentration of $\mu$ into many close prefix sums with
opposite first bits.  Opposite first bits ensure that their difference
polynomial has constant coefficient $-1$, as required by the algebraic
criterion.

\subsection{Overlap of the first-bit measures}

Define subprobability measures
\[
  \mu_i(A)=\mathbb P(X\in A,\ u_0=i),\qquad i=0,1.
\]
Their sum is $\mu$.  The following fact uses only $\lambda>1/2$, not
singularity.

\begin{lemma}\label{cls:first-bit-overlap}
There is a finite measure $\xi\ne0$ with $\xi\leq\mu_0$ and
$\xi\leq\mu_1$.
\end{lemma}

\begin{proof}
If no such measure exists, the densities of $\mu_0$ and $\mu_1$ with
respect to $\mu$ have disjoint supports almost everywhere.  Thus
$\mu_0$ and $\mu_1$ are mutually singular.  We will show that this
makes the conditional entropy of the first $n$ bits, given a
$\lambda^n$-quantization of $X$, arbitrarily small compared with $n$.
There are only $O(\lambda^{-n})$ quantization cells, too few to carry
the full entropy $n\log2$ of those bits.
For every $\eta\in(0,1/2)$,
inner regularity then gives disjoint compact sets $K_0,K_1$ such that
\[
  \mathbb P(X\in K_{u_0})>1-\eta.
\]
Let $\delta=\operatorname{dist}(K_0,K_1)>0$, and choose an integer
$h\geq1$ with $\lambda^h<\delta/3$.

Set $Q_n=\lfloor X/\lambda^n\rfloor$.  Given $Q_n$ and the bits
$u_0,\ldots,u_{i-1}$, one can approximate
\[
  X^{(i)}=\frac{X-\sum_{k<i}u_k\lambda^k}{\lambda^i}
\]
to within $\lambda^{n-i}$.  The pair $(X^{(i)},u_i)$ has the same law
as $(X,u_0)$.  For $i\leq n-h$, the approximation therefore determines
an estimator of $u_i$ with error probability at most $\eta$: when
$X^{(i)}\in K_{u_i}$, choose the nearer of the two compact sets.

Write $H$ for Shannon entropy using natural logarithms, and put
$h_2(\eta)=-\eta\log\eta-(1-\eta)\log(1-\eta)$.
The error-indicator argument gives
\[
  H(u_i\mid Q_n,u_0,\ldots,u_{i-1})\leq h_2(\eta)
  \qquad(i\leq n-h).
\]
Indeed, the estimator and its binary error indicator recover $u_i$;
the entropy of that indicator is at most $h_2(\eta)$.
The remaining at most $h$ bits have conditional entropy at most
$h\log2$.  Since $X\in[0,L]$, the variable $Q_n$ takes at most
$L\lambda^{-n}+2$ values.  The chain rule now gives
\[
  n\log2
  \leq H(Q_n)+H(u_0,\ldots,u_{n-1}\mid Q_n)
  \leq n\log(1/\lambda)+n h_2(\eta)+O_{\lambda,\eta}(1).
\]
Choose $\eta$ so small that $h_2(\eta)<\log(2\lambda)$ and let
$n\to\infty$.  This is a contradiction.
\end{proof}

\begin{proposition}[Many close prefixes with opposite first bits]
\label{cls:cross-pairs}
If $\mu_\lambda$ is singular, there are constants $a\in(0,L/2)$ and
$c>0$, depending only on $\lambda$, with the following property.
For every $Q>1$ and all sufficiently large $n$, at least $c2^n$ words
$u\in\{0,1\}^n$ with $u_0=0$ each have at least $Q(2\lambda)^n$
distinct neighbors $v\in\{0,1\}^n$ such that
\begin{equation}
  v_0=1,
  \qquad |s_u-s_v|\leq(L-a)\lambda^n.
  \label{cls:cross-pair-bound}
\end{equation}
The constants $a,c$ are independent of $Q$ and $n$.
\end{proposition}

\begin{proof}
Fix $\xi$ from Lemma~\ref{cls:first-bit-overlap}, and write
$b=\xi(\mathbb R)>0$.  The endpoints $0,L$ have $\mu$-measure zero,
because each requires all the bits to have one prescribed value.
Choose $a\in(0,L/2)$ so that
$\mu([a,L-a])>1-b/8$.

For each $n$, let
$X^{(n)}=\sum_{k\geq0}u_{n+k}\lambda^k$, and restrict the first-bit
measures by defining
\[
  \mu'_{i,n}(A)=
  \mathbb P(X\in A,\ u_0=i,\ X^{(n)}\in[a,L-a]).
\]
The sum of the masses discarded from $\mu_0,\mu_1$ is less than
$b/8$, since $X^{(n)}$ has law $\mu$.

Fix $Q>1$.  By Lemma~\ref{cls:compact-concentration}, choose a compact
set $K$ with $\operatorname{Leb}(K)<\varepsilon$ and
$\mu(K)>1-\varepsilon$, where
\[
  \varepsilon<b/8,
  \qquad \frac{2\varepsilon Q}{a}<\frac b4.
\]
Then $\xi(K)>7b/8$.  Partition the line into half-open intervals of
length $a\lambda^n$, and let $\mathcal J_n$ consist of those meeting
$K$.  Compactness of $K$ implies, for all sufficiently large $n$,
\[
  \#\mathcal J_n\leq\frac{2\varepsilon}{a\lambda^n}.
\]
To see this, their union is contained in the closed
$a\lambda^n$-neighborhood of $K$, whose Lebesgue measure tends to
$\operatorname{Leb}(K)$.  Moreover,
\begin{equation}
  \sum_{J\in\mathcal J_n}
       \min_{i=0,1}\mu'_{i,n}(J)\geq\frac{3b}{4}.
  \label{cls:common-bin-mass}
\end{equation}
Indeed, the corresponding sum with $\mu_i$ in place of $\mu'_{i,n}$
is at least $\xi(K)$, and restricting the two measures loses less
than $b/8$ in total.

Let $U_i(J)$ be the set of words $u$ of length $n$ with $u_0=i$ for
which
\[
  \bigl(s_u+\lambda^n[a,L-a]\bigr)\cap J\ne\varnothing.
\]
Each prefix has probability $2^{-n}$, so
$\mu'_{i,n}(J)\leq2^{-n}\#U_i(J)$.  Discard the intervals for which
$\#U_1(J)<Q(2\lambda)^n$.  Their contribution to
\eqref{cls:common-bin-mass} is at most
\[
  \frac{2\varepsilon}{a\lambda^n}\,
  2^{-n}Q(2\lambda)^n
  =\frac{2\varepsilon Q}{a}<\frac b4.
\]
Thus the sum of the minima over the remaining intervals is at least
$b/2$.

A fixed interval $s_u+\lambda^n[a,L-a]$ meets at most
\[
  C_a=\left\lceil\frac{L-2a}{a}\right\rceil+2
\]
of the bins.  If $\mathcal U$ is the union of $U_0(J)$ over the retained
bins, it follows that
\[
  \frac b2\leq2^{-n}\sum_{J\text{ retained}}\#U_0(J)
       \leq2^{-n}C_a\#\mathcal U.
\]
Take $c=b/(2C_a)$.  For every $u\in\mathcal U$, choose one retained
bin $J$ with $u\in U_0(J)$.  All $v\in U_1(J)$ satisfy
\eqref{cls:cross-pair-bound}: two points in $J$ differ by less than
$a\lambda^n$, and two tails in $\lambda^n[a,L-a]$ differ by at most
$(L-2a)\lambda^n$.  There are at least $Q(2\lambda)^n$ such neighbors.
\end{proof}

\subsection{Completion to irreducible polynomials}
\label{cls:subsection-completion}

For every pair in Proposition~\ref{cls:cross-pairs}, the difference
$w=u-v$ has $w_0=-1$ and
$|\sum_{i<n}w_i\lambda^i|\leq(L-a)\lambda^n$.
We next extend any such coefficient vector to a polynomial whose root
lies in $(\lambda,\lambda+4^{-n}/2]$.  The first lemma ensures that a
bounded-coefficient polynomial cannot remain too flat throughout this
short interval.

\begin{lemma}[Uniform finite-order nonvanishing]
\label{cls:finite-order}
Fix $\lambda\in(0,1)$.  There are an integer $M\geq1$ and constants
$c_0,x_0>0$ such that every real polynomial
$p(t)=-1+\sum_{k\geq1}a_kt^k$ with $|a_k|\leq1$ satisfies
\[
  \sup_{t\in[\lambda,\lambda+x]}|p(t)|\geq c_0x^M
  \qquad(0<x<x_0).
\]
\end{lemma}

\begin{proof}
There exist $M\geq1$ and $\delta>0$ such that
\begin{equation}
  \max_{0\leq k\leq M}
     \left|\frac{p^{(k)}(\lambda)}{k!}\right|\geq\delta
  \label{cls:jet-bound}
\end{equation}
for every polynomial in the stated class.  Otherwise, for each $m$ one
could choose a polynomial for which all these quantities through order
$m$ are less than $1/m$.  Extend its coefficient sequence by zeros.
A coefficientwise convergent subsequence exists by compactness of
$[-1,1]^{\mathbb N}$.  The associated series converge, with every
derivative, uniformly on compact subsets of the unit disk.  Their limit
would have all derivatives zero at $\lambda$ and constant coefficient
$-1$, contradicting the identity theorem.

For fixed $M$, equivalence of norms on polynomials of degree at most
$M$ gives $\kappa_M>0$ such that
\[
  \sup_{0\leq t\leq1}\left|\sum_{k=0}^M b_kt^k\right|
       \geq\kappa_M\max_{k\leq M}|b_k|.
\]
Apply this to the Taylor polynomial of $p(\lambda+xt)$.
For $x\leq1$, \eqref{cls:jet-bound} makes the right-hand side at least
$\kappa_M\delta x^M$.  The coefficient bound gives a uniform bound on
the $(M+1)$st derivative in a fixed neighborhood of $\lambda$, so the
Taylor remainder is $O_\lambda(x^{M+1})$, with $M$ fixed.  Reducing
$x_0$ proves the assertion with $c_0=\kappa_M\delta/2$.
\end{proof}

To impose irreducibility without disturbing the prescribed coefficients,
we use the prime-polynomial theorem in arithmetic progressions.  We
include the deduction needed here from the function-field Riemann
hypothesis; its character-sum formulation is recalled in
\cite[Theorem~3.2 and Lemma~3.3]{GorodetskyKovaleva2024}.

\begin{lemma}[Irreducibles with prescribed initial coefficients]
\label{cls:prescribed-prime}
Fix integers $D\geq2$ and $d>2D$.  For all sufficiently large $n$,
every polynomial $F\in\mathbb F_2[t]$ with $\deg F<Dn$ and $F(0)=1$
has a monic irreducible completion $P\in\mathbb F_2[t]$ satisfying
\[
  \deg P=dn,
  \qquad P\equiv F+t^{Dn}\pmod{t^{Dn+1}}.
\]
The threshold in $n$ is uniform in $F$.
\end{lemma}

\begin{proof}
For a monic polynomial $f$, write $\Lambda(f)=\deg P$ if $f=P^k$
for a monic irreducible $P$, and $\Lambda(f)=0$ otherwise.
Let $\chi$ be a nonprincipal Dirichlet character modulo $t^\ell$,
extended by zero to nonunits.  Its $L$-function is a polynomial of
degree at most $\ell-1$, with inverse roots $\alpha_j$ satisfying
$|\alpha_j|\leq\sqrt2$; this bound includes trivial factors.
Taking the logarithmic derivative of the Euler product gives
\[
  \left|\sum_{\substack{f\text{ monic}\\\deg f=s}}
       \Lambda(f)\chi(f)\right|
  =\left|\sum_j\alpha_j^s\right|
  \leq\ell2^{s/2}.
\]
For the principal character, the same sum is $2^s-1$, by the zeta
function of $\mathbb F_2[t]$ after removal of the prime $t$.
Character orthogonality therefore gives, in every invertible residue
class $A$ modulo $t^\ell$,
\begin{equation}
  \sum_{\substack{f\text{ monic},\ \deg f=s\\f\equiv A\bmod t^\ell}}
       \Lambda(f)
  \geq \frac{2^s-1}{2^{\ell-1}}-\ell2^{s/2}.
  \label{cls:prime-progression-bound}
\end{equation}
The contribution from proper prime powers, even without the congruence
restriction, is at most
\[
  \sum_{\substack{k\mid s\\k\geq2}}2^{s/k}
       \leq s2^{s/2}.
\]
Take $\ell=Dn+1$, $s=dn$, and $A=F+t^{Dn}$.  This is an invertible
class.  The main term in \eqref{cls:prime-progression-bound} has order
$2^{(d-D)n}$, whereas the character error and proper-power contribution
are $O(n2^{dn/2})$.  Since $d>2D$, at least one irreducible polynomial
remains for all sufficiently large $n$, uniformly in $A$.
\end{proof}

\begin{proposition}[Realization of a polynomial prefix]
\label{cls:realization}
Fix $\lambda\in(1/2,1)$ and $a\in(0,L)$, where $L=(1-\lambda)^{-1}$.
There is an integer $d_0$ such that, for each integer $d\geq d_0$ and
all sufficiently large $n$, the following assertion holds uniformly
over $w\in\{-1,0,1\}^n$.  If
\begin{equation}
  w_0=-1,
  \qquad
  \left|\sum_{i=0}^{n-1}w_i\lambda^i\right|
      \leq(L-a)\lambda^n,
  \label{cls:realization-hypothesis}
\end{equation}
then some $r\in R(n,d)$ satisfies
\[
  \lambda<r\leq\lambda+\tfrac12 4^{-n},
  \qquad \pi_n(p_r)=w.
\]
\end{proposition}

\begin{proof}
We first append digits to make the polynomial small at $\lambda$.
We then construct two irreducible completions with opposite signs there
but the same sign at a nearby point to its right.  Uniform finite-order
nonvanishing supplies that second point, and the intermediate value
theorem supplies a root between the two points.

Let $M,c_0,x_0$ be supplied by Lemma~\ref{cls:finite-order}.  Choose
an integer $D\geq2$ with
\begin{equation}
  \lambda^D<4^{-M},
  \label{cls:D-choice}
\end{equation}
and then take $d_0=2D+1$.

Start with the polynomial in \eqref{cls:realization-hypothesis} and
append coefficients through index $Dn-1$.  If the next index is $k$,
let $e$ be the negative of the current value at $\lambda$, divided by
$\lambda^k$.  Appending $e_0\in\{-1,0,1\}$ changes this residual to
\[
  e\longmapsto\frac{e-e_0}{\lambda}.
\]
Choose a nearest digit $e_0$.  Once $|e|\leq1$, the next residual has
absolute value at most $1/(2\lambda)<1$, so this bound is preserved.
Before that time, the residual keeps its sign and its absolute value
updates by
\[
  |e|\longmapsto\frac{|e|-1}{\lambda}
        =L+\frac{|e|-L}{\lambda}.
\]
The initial bound $|e|\leq L-a$ implies that $|e|\leq1$ is reached
within a number of steps depending only on $\lambda,a$.  Indeed,
after $k$ steps without reaching this interval the absolute value is
at most $L-a\lambda^{-k}$, which eventually falls below $1$.
Consequently, for all sufficiently large $n$, the polynomial $p_0$
constructed through index $Dn-1$ satisfies
\begin{equation}
  |p_0(\lambda)|\leq\frac{\lambda^{Dn}}{2\lambda}.
  \label{cls:residual}
\end{equation}

Put $x_n=4^{-n}/2$.  By Lemma~\ref{cls:finite-order}, there is
$t_n\in[\lambda,\lambda+x_n]$ with
$|p_0(t_n)|\geq c_0x_n^M$.  Meanwhile every possible appendage with
coefficients of absolute value at most one, starting at index $Dn$,
is bounded throughout this interval by
\[
  \frac{(\lambda+x_n)^{Dn}}{1-\lambda-x_n}=o(x_n^M).
\]
This follows from \eqref{cls:D-choice}, since
$Dn x_n\to0$.  The estimate is uniform in the prefix and in the
eventual length of the appendage.

Fix $d\geq d_0$.  Apply Lemma~\ref{cls:prescribed-prime} to the
reduction of $p_0$ modulo $2$, obtaining a monic irreducible polynomial
$P$ of degree $dn$ with
$P\equiv p_0+t^{Dn}\pmod{t^{Dn+1}}$.
Represent its coefficients of indices $Dn,\ldots,dn$ by integers
$b_k\in\{0,1\}$, and let
\[
  T(t)=\sum_{k=Dn}^{dn}b_kt^k,
  \qquad p_\pm(t)=p_0(t)\pm T(t).
\]
Here $b_{Dn}=b_{dn}=1$.  Thus $T(\lambda)\geq\lambda^{Dn}$, and
\eqref{cls:residual} implies
$p_-(\lambda)<0<p_+(\lambda)$.  At $t_n$, the negligible-tail
estimate shows that both polynomials have the same nonzero sign as
$p_0(t_n)$.  One of them therefore has a real root
\[
  r\in(\lambda,t_n]\subset(\lambda,\lambda+x_n].
\]

Both $p_\pm$ reduce to $P$ modulo $2$, so they are irreducible over
$\mathbb Q$.  They have degree $dn$, leading coefficient $\pm1$,
constant coefficient $-1$, and all coefficients in $\{-1,0,1\}$.
Their roots are algebraic units.  For sufficiently large $n$, the
chosen $r$ is in $(1/2,1)$, belongs to $R(n,d)$, and its signed minimal
polynomial is the chosen $p_\pm$.  Its first $n$ coefficients are $w$.
\end{proof}

\subsection{Proof of the converse and the quantifiers}

Suppose $\mu_\lambda$ is singular.  Proposition~\ref{cls:cross-pairs}
gives constants $a,c>0$.  Choose one integer $d$ large enough for
Proposition~\ref{cls:realization} and also satisfying $1/d<c$.
This choice depends on $\lambda$ alone.

Fix any integer $j\geq1$ and choose $Q>2j$.  For all sufficiently
large $n$, each of at least $c2^n$ words $u$ has at least
$Q(2\lambda)^n$ neighbors $v$ supplied by
Proposition~\ref{cls:cross-pairs}.  For each difference $w=u-v$,
Proposition~\ref{cls:realization} supplies a root $r_w\in R(n,d)$
with signed minimal polynomial projecting to $w$, and
\[
  \lambda<r_w\leq\lambda+\tfrac12 4^{-n}.
\]
Choose any one of these roots as a center $r$.  All of them lie within
$4^{-n}$ of $r$, so all their projected vectors belong to $W_n(r;d)$.
For fixed $u$, its different neighbors give distinct vectors $u-v$,
each in $\mathcal O_u$.  Furthermore,
\[
  1\leq\left(\frac r\lambda\right)^n
       \leq\left(1+\frac{4^{-n}}{2\lambda}\right)^n\longrightarrow1
\]
uniformly in the choice of center.  Hence, for all sufficiently large
$n$, each of these $c2^n$ orthants contains at least
$j(2r)^n$ distinct vectors of $W_n(r;d)$.  Since $c>1/d$, we have
$r\in R_*(n,d,j)$ and
$\lambda\in[r-4^{-n},r)$.

We have proved that a single $d$ works for every $j$, at every
sufficiently large $n$.  This proves membership in the expression
\eqref{cls:criterion} with $\liminf$.  Membership in that expression
implies membership with $\limsup$, which implies singularity by
Section~\ref{cls:subsection-forward}.  Lemma~\ref{cls:pure-type}
handles $\lambda\leq1/2$ and supplies the absolute continuity of every
remaining parameter.  This completes the proof of
Theorem~\ref{cls:classification}.

\section{Recovery of classical parameter classes}
\label{cls:section-classical}

The arithmetic criterion is compatible with the classical norm
separation argument for Garsia parameters and the Fourier argument for
Pisot parameters.  The first example directly bounds the crowding in
\eqref{cls:retained-roots}; the second establishes singularity and then
uses Theorem~\ref{cls:classification} to obtain the approximation
property.

\begin{proposition}[Garsia parameters, \cite{Garsia1962}]
\label{cls:garsia}
Let $\beta\in(1,2)$ be an algebraic integer of absolute norm $2$ whose
conjugates all have modulus greater than $1$.  Then $1/\beta\notin
\mathcal C$, and $\nu_{1/\beta}$ is absolutely continuous.
\end{proposition}

\begin{proof}
Put $\lambda=1/\beta$.  Consider distinct vectors
$w,w'\in\{-1,0,1\}^n$ lying in one closed orthant.  Their difference
has coefficients in $\{-1,0,1\}$, so
\[
  P(t)=\sum_{i=0}^{n-1}(w_i-w'_i)t^{n-1-i}
\]
is a nonzero polynomial of that coefficient type.  It cannot vanish
at $\beta$.  Otherwise, after removing any power of $t$ dividing
$P$, its constant coefficient would be $\pm1$, while divisibility by
the monic minimal polynomial of $\beta$ would force that coefficient
to be divisible by $2$.

For every other conjugate $\gamma$ of $\beta$,
\[
  |P(\gamma)|\leq\frac{|\gamma|^n}{|\gamma|-1}.
\]
The nonzero algebraic integer $P(\beta)$ has absolute norm at least
one.  Since the product of the moduli of all conjugates of $\beta$
is $2$, the preceding estimate gives a constant $c_\beta>0$ such that
\begin{equation}
  \left|\sum_{i=0}^{n-1}(w_i-w'_i)\lambda^i\right|
  =\beta^{-(n-1)}|P(\beta)|\geq c_\beta2^{-n}.
  \label{cls:garsia-spacing}
\end{equation}

For any potential hit of \eqref{cls:criterion} at a sufficiently large
$n$, all projected vectors satisfy the slab bound \eqref{cls:slab}.
Within any one orthant, \eqref{cls:garsia-spacing} bounds their number
by
\[
  1+\frac{2B\lambda^n}{c_\beta2^{-n}}
       \leq C_\beta(2\lambda)^n.
\]
Choose an integer $j>C_\beta$.  Since a hit has $r>\lambda$, no
orthant can then meet the threshold $j(2r)^n$.  This excludes all
sufficiently large $n$, for every degree multiplier $d$.  Thus
$\lambda\notin\mathcal C$, and Theorem~\ref{cls:classification}
gives absolute continuity.
\end{proof}

\begin{proposition}[Reciprocal Pisot parameters, \cite{Erdos1939}]
\label{cls:pisot}
If $\beta\in(1,2)$ is a Pisot number, then $1/\beta\in\mathcal C$,
and $\nu_{1/\beta}$ is singular.
\end{proposition}

\begin{proof}
A Pisot number is an algebraic integer greater than one whose other
conjugates have modulus less than one.  Here its nonzero integer norm
has absolute value less than $\beta<2$, so $\beta$ is a unit.
Consequently $\beta^\ell$ is an algebraic integer for every
$\ell\in\mathbb Z$, and it can never be a half-integer of odd
numerator.  In particular, all factors $\cos(\pi\beta^\ell)$ are
nonzero.

For negative $\ell$, the quantity $\beta^\ell$ tends exponentially
to zero.  For positive $\ell$, its distance to an integer tends
exponentially to zero: the trace of $\beta^\ell$ is an integer and
the contributions of the other conjugates decay exponentially.
The quadratic behavior of $-\log|\cos(\pi t)|$ near the integers
therefore gives
\[
  \sum_{\ell\in\mathbb Z}-\log|\cos(\pi\beta^\ell)|<\infty.
\]
With $\lambda=1/\beta$, independence of the signs yields
\[
  \left|\mathbb E\exp(i\pi\beta^nY_\lambda)\right|
     =\prod_{\ell\leq n}|\cos(\pi\beta^\ell)|
     \geq\prod_{\ell\in\mathbb Z}|\cos(\pi\beta^\ell)|>0.
\]
The Riemann--Lebesgue lemma rules out absolute continuity.
Lemma~\ref{cls:pure-type} gives singularity, and
Theorem~\ref{cls:classification} gives $\lambda\in\mathcal C$.
\end{proof}
